\documentclass[12pt]{article}
\usepackage[T1]{fontenc}
\usepackage[utf8]{inputenc}
\usepackage{lmodern}
\usepackage{textcomp}
\usepackage{enumitem}
\usepackage{geometry}
\geometry{margin=1in}
\begin{document}
\begin{center}
Answer Key - History - Undergraduate Level Assessment 19 \
\small (Answers are ordered exactly as in the question paper.)
\end{center}

\section*{Multiple Choice}
\begin{enumerate}[label=\arabic*.]
\item \textbf{Answer:} D
\\ \textit{Options:} A) pottery fragments.; B) irrigation canals.; C) jewelry artifacts.; D) monumental works.; E) pyramids.
\\ \textit{Note:} Monumental works, such as temples, palaces, and fortifications, serve as tangible representations of the power, resources, and cultural values of ancient state societies, making them the most obvious material symbols of these complex civilizations.
\item \textbf{Answer:} A
\\ \textit{Options:} A) was driven solely by population growth.; B) only occurred in Europe and Asia.; C) took hundreds of thousands of years.; D) always required careful planning.; E) happened everywhere except in South America.
\\ \textit{Note:} The transition from food gathering to total reliance on agriculture was primarily driven by population growth, as increasing numbers of people created a greater demand for stable food sources, prompting societies to develop agricultural practices for sustained food production.
\item \textbf{Answer:} E
\\ \textit{Options:} A) Erectus was not bipedal.; B) Erectus was not capable of using tools.; C) Erectus was primarily herbivorous.; D) Erectus fossils are found only in Africa.; E) Erectus possessed a larger brain.
\\ \textit{Note:} The correct answer is supported by the fact that Homo erectus had a significantly larger brain size, averaging around 900 to 1,100 cubic centimeters, compared to Homo habilis, which had a brain size averaging about 510 to 600 cubic centimeters, indicating a more advanced cognitive ability and adaptation.
\item \textbf{Answer:} D
\\ \textit{Options:} A) both a and b; B) Clovis; C) Gravettian; D) Mousterian; E) Acheulean
\\ \textit{Note:} The Mousterian tool technology, characterized by a sophisticated flake-based manufacturing technique, is specifically associated with Neandertals, reflecting their advanced cognitive abilities and adaptation to diverse environments during the Middle Paleolithic period.
\item \textbf{Answer:} E
\\ \textit{Options:} A) a shared ability to use complex tools; B) a similar diet; C) a reliance on brachiation; D) an identical social structure; E) a mode of locomotion
\\ \textit{Note:} Homo sapiens and Australopithecus afarensis both share a bipedal mode of locomotion, which is a significant characteristic of hominins that distinguishes them from other primates.
\end{enumerate}

\section*{True / False}
\begin{enumerate}[label=\arabic*.]
\setcounter{enumi}{5}
\item \textbf{Answer:} False
\\ \textit{Options:} A) True; B) False
\\ \textit{Note:} The correct answer is: 1332
\item \textbf{Answer:} False
\\ \textit{Options:} A) True; B) False
\\ \textit{Note:} The correct answer is: Tsongkhapa
\item \textbf{Answer:} True
\\ \textit{Options:} A) True; B) False
\\ \textit{Note:} According to the passage: an apology
\item \textbf{Answer:} False
\\ \textit{Options:} A) True; B) False
\\ \textit{Note:} The correct answer is: Zaynab
\item \textbf{Answer:} True
\\ \textit{Options:} A) True; B) False
\\ \textit{Note:} According to the passage: The National Historical Publications and Records Commission
\end{enumerate}

\section*{Long Form Response}
\begin{enumerate}[label=\arabic*.]
\setcounter{enumi}{10}
\item \textbf{Answer:} Iberian peninsula
\\ \textit{Note:} Expected answer: Iberian peninsula
\item \textbf{Answer:} quoting Plato
\\ \textit{Note:} Expected answer: quoting Plato
\end{enumerate}

\vfill
\noindent\textit{End of Answer Key}
\end{document}